\section{Related Work and Background}
\label{sec:related_work}


\subsection{Hyperdimensional Computing (HDC)}
\label{subsec:hdc_background}

Hyperdimensional computing (HDC) is a brain-inspired computational paradigm in which symbols and structured entities are represented as high-dimensional vectors---called \emph{hypervectors}---with components drawn independently from simple distributions~\cite{kanerva2009hyperdimensional}. In such spaces, independently sampled hypervectors are nearly orthogonal with high probability, a geometric property that underlies classical HDC operations such as \emph{bundling} (superposition via addition), \emph{binding} (association via elementwise multiplication or permutation), and \emph{permutation} (role shifting). Because information is distributed holographically across all dimensions, HDC representations exhibit strong robustness to noise, quantization, and partial corruption.

Given an input $x$, an HDC encoder produces a bounded hypervector $h=\phi(x)\in\mathbb{R}^D$ (or $\mathbb{C}^D$) using random projections or compositional schemes built from a small set of base hypervectors. These encoders admit a precise random-feature interpretation: the empirical kernel
\begin{equation}
\label{eq:empirical_kernel}
k_D(x,x')=\langle \phi(x),\phi(x')\rangle
\end{equation}
converges to a smooth limiting kernel $k_\ast(x,x')$ as $D\to\infty$, with uniform concentration at rate $O(D^{-1/2})$.

The following concentration property, standard in the random-feature literature, quantifies this convergence.

With probability at least $1-\delta$ over the random generation of hypervectors, the empirical kernel concentrates uniformly around the limiting kernel:
\begin{equation}
\label{eq:kernel_concentration}
\big| k_D(x,x') - k^\ast(x,x') \big| \leq \frac{C_k}{\sqrt{D}},
\quad \forall\, x,x' \in \mathcal{S} \times \mathcal{A},
\end{equation}
where $C_k > 0$ is an absolute constant depending only on the encoding scheme.

This $O(D^{-1/2})$ concentration rate is established for various random-feature constructions, including Random Fourier Features and related schemes~\cite{rahimi2007random,bach2015equivalence,rudi2017generalization}. As a consequence, linear prediction in the hypervector domain serves as a computationally efficient, finite-dimensional approximation to kernel methods in the RKHS associated with $k_\ast$.




\cite{ni2022qhd} proposes a hyperdimensional regression model to approximate the Q-value function.

\cite{ye2023fedfm} anchor based feature matching FRL

\cite{jin2022federated}  Federated Reinforcement Learning with Environment Heterogeneity (FRL-EH) framework

\cite{hwang2025fed} Federated Reinforcement Learning with Environment Heterogeneity (FRL-EH) framework by incorporating a regularization term in the local update